% Generated from paper/source. Do not edit by hand.
\section{Mathematical framework}
\subsection{Complex hypothesis-state update}
Let a tokenized observation at step t be embedded as a real vector x\_t. The implemented complex operator maintains h\_t in C\textasciicircum{}d and applies a token-conditioned update followed by a phase-compatible nonlinearity and a magnitude-sensitive readout. Suppressing low-rank factorization and bias terms, the linear part is

\begin{equation}
h_t = W_t h_{t-1} + U x_t
\end{equation}

Write h = a + ib and W = A + iB, with a, b real vectors and A, B real matrices. Then the same multiplication is

\begin{equation}
\begin{bmatrix}\Re(Wh)\\\Im(Wh)\end{bmatrix}=\begin{bmatrix}A&-B\\B&A\end{bmatrix}\begin{bmatrix}a\\b\end{bmatrix}
\end{equation}

This is an identity, not an approximation. The exact real block has the same real degrees of freedom as W because the four blocks are tied by A and B. When initialization, token order, loss, and optimizer policy are mapped consistently, any prediction difference beyond numerical tolerance is an implementation failure or a consequence of a non-mapped operation, not evidence for uniquely complex capacity.

\subsection{Three distinct hypotheses}
The comparator audit separates three hypotheses that are often conflated.

\begin{itemize}
\item \textbf{Representational uniqueness:} the complex model computes a function unavailable to the exact structured real block.
\item \textbf{Parameterization advantage:} the complex coordinates create a more useful inductive bias or optimization trajectory even though the represented mapped function class overlaps.
\item \textbf{Architecture advantage:} the complete complex operator outperforms all parameter- and compute-matched real sequence architectures on the prespecified outcome.
\end{itemize}

The algebra directly falsifies representational uniqueness for the mapped linear complex operations. Empirical mapped runs test whether the implementation preserves that identity through the complete network. The best-real envelope tests architecture advantage over the real models actually evaluated. Parameterization advantage remains a legitimate possibility in other objectives, initializations, nonlinearities, or training regimes; it is not supported by the present top-1 results.

\subsection{Primary estimand}
For family f, training size n, training seed s, world w, and severity q, let A\_C(f,n,s,w,q) be complex top-1 accuracy. Let R be the prespecified set of evaluated real controls. The nested effect is

\begin{equation}
\Delta_{fnswq}=A_C(f,n,s,w,q)-\max_{r\in\mathcal{R}}A_r(f,n,s,w,q)
\end{equation}

This cellwise envelope is intentionally adversarial. It asks whether the candidate beats the strongest real alternative in each matched condition, not whether it beats a conveniently chosen average control. It also creates selection optimism for the real side, which is acceptable for a falsifier but must be acknowledged when interpreting magnitude.

\subsection{Structural predictors and candidate law}
Each generator world exposes two preregistered structural descriptors: empirical mutual information between observed order and the target, and shift severity q. Six candidate families were compared on discovery aggregates. The selected quadratic candidate was

\begin{equation}
\widehat{\Delta}=\beta_0+\beta_1 I+\beta_2 q+\beta_3 I^2+\beta_4 Iq+\beta_5 q^2
\end{equation}

The coefficients were frozen before held-out family evaluation. Passing required R2 at least 0.50, sign accuracy at least 0.80, and mean absolute error at most 0.015. A fit to discovery data was not itself treated as a law.

\subsection{Claim rule}
The preregistered practical threshold for a favorable complex effect was +0.02 top-1. Because the executed independent profiles deviate from the full protocol, they are outcome-ineligible for the strongest confirmatory category even if a threshold were crossed. Exact equivalence, empirical effect direction, and law failure are reported separately rather than collapsed into one verdict.
